\section{Caso de los preesquemas no conexos: categorías multigaloisianas}
\label{V.9}

Sea~$S$ un preesquema localmente noetheriano, y sean
$(S_i)_{i\in I}$ sus componentes conexas. Entonces la categoría
$\cal{C}(S)$ de los recubrimientos étales de~$S$ es equivalente
al producto de las categorías~$\cal{C}(S_i)$, que se interpretan
en términos de los grupos fundamentales de las~$S_i$, una vez
\ifthenelse{\boolean{orig}}{elegido}{elegido} un punto geométrico en
cada~$S_i$. En la aplicación de la teoría del descenso a los morfismos
étales, resulta a veces incómodo elegir un punto geométrico de~$S_i$
para cada~$S_i$. Es entonces más cómodo recurrir a la generalización
natural de~\Ref{V.5.8} para interpretar~$\cal{C}(S)$ como una categoría
de funtores sobre el grupoide de los puntos geométricos de~$S$,
considerado como suma de los grupoides correspondientes a las componentes
conexas de~$S$; los funtores en cuestión son los funtores con valores en
la categoría de los conjuntos finitos que satisfacen la propiedad de
continuidad análoga a la invocada en~\Ref{V.5.8}. En la práctica, se tendrá
una familia~$(a_t)_{t\in E}$ de puntos geométricos de~$S$, tal que toda
componente conexa~$S_i$ de~$S$ contenga al menos uno, y podremos entonces,
como en~\Ref{V.5.10}, reemplazar el grupoide de todos los puntos geométricos
de~$S$ por el grupoide análogo cuyo conjunto subyacente es~$E$. Por supuesto,
estas consideraciones deberían insertarse en definiciones generales acerca
de las categorías equivalentes a productos de categorías de la
forma~$\cal{C}(\pi)$, que podrán llamarse
\emph{categorías multigaloisianas}.
\index{multigaloisiana (categoría)|hyperpage}%
Dejaremos los detalles al lector.
